\guidepart{teal}{Connect}{Connect to all the things}
  {Eight ways in. Pick the one you need; skip the rest.}
  {\partitem{\iTerminal}{SSH: the big one}
   \partitem{\iUsb}{Serial: cables and consoles}
   \partitem{\iPlug}{Telnet and RAW TCP}
   \partitem{\iDesktop}{VNC remote desktops}
   \partitem{\iChat}{IRC}
   \partitem{\iGithub}{GitHub}}

% =========================================================================================
\chapter{SSH: the big one}\label{ch:ssh}
\chapmeta{10 minutes (or 2 for the basics)}{You'll log in with keys, hop through bastions and never type a 2FA code again}

\begin{tldr}
Host, username, key, done. Everything else (jump hosts, proxies, forwards, 2FA, login scripts) is optional and
lives in its own fold-away section of the editor. Open only the ones you need.
\end{tldr}

\begin{glance}{teal}{SSH}
  Default port & 22\\
  Logs in with & Password, public key, keyboard-interactive, SSH agent\\
  Extras       & SFTP file browser, port forwards, live latency, agent forwarding, TOTP 2FA\\
  Good for     & Servers, VMs, Raspberry Pis, anything with \code{sshd}\\
\end{glance}

\section{The two-minute version}

\begin{enumerate}
  \item \keys{\mod+N}, leave the protocol on \proto[teal]{SSH}.
  \item \ui{Host}, \ui{Username}. Leave the port on 22 unless you know otherwise.
  \item Under \ui{Key \& Passphrase}, pick your key from the suggestions (Nexus Terminal lists the keys in your
        \code{\textasciitilde/.ssh} folder) and check for a little green \ui{Key verified}.
  \item \ui{Create \& Connect}. The first time, you'll be asked to trust the server's host key. Say yes if you
        were expecting it.
\end{enumerate}

\crop{hostkey.png}{842}{720}{845}{589}{0.66\linewidth}{First time connecting? You'll be shown the server's fingerprint. \ui{Accept \& Save} if it's the server you meant.}

\begin{headsup}
Type the full path to a key if you're not using the suggestions. Nexus Terminal doesn't expand \code{\textasciitilde}
in that box, so write \code{/home/you/.ssh/id\_ed25519} (or \code{C:\textbackslash Users\textbackslash you\textbackslash .ssh\textbackslash id\_ed25519}).
\end{headsup}

\section{How it logs in}

Under \ui{Authentication}, tick which \ui{Auth Methods} to try: \ui{Public key}, \ui{Password} and
\ui{Keyboard-interactive}. Under \ui{SSH Agent}, tick \ui{Use SSH Agent for Authentication} to use keys already
loaded into \code{ssh-agent} (or Pageant, or the Windows OpenSSH agent).

Every password and passphrase has a \emph{mode}: save it securely, remember it until you quit, or ask every time.
That's all explained in \hyperref[ch:secrets]{Passwords \& the vault} (p.\,\pageref{ch:secrets}).

\begin{protip}[Agent forwarding]
Tick \ui{Forward Agent to Remote} (it's \code{ssh -A}) and you can \code{git pull} or hop onwards from the server using
the keys on \emph{your} machine. Only do it on servers you trust: anyone with root there can use your agent while
you're connected.
\end{protip}

\section{Jump hosts: getting through the bastion}

Lots of servers hide behind a bastion. Add one or more hops under \ui{Jump Hosts}, each with its own user, host,
port and optional key. Nexus Terminal tunnels through them in order, so the chain
\textit{you\then bastion\then internal-box\then db-primary} is just three rows.

\section{Proxies}

Stuck behind a corporate proxy? Under \ui{Proxy}, choose \ui{SOCKS5} or \ui{HTTP CONNECT} and fill in the host,
port and (optional) username and password. The proxy is used to reach the SSH server, or the first jump host if
you have some.

\section{2FA without the phone faff}\label{sec:totp}

If your server asks for a one-time code, paste the Base32 TOTP secret (the text version of the QR code) into
\ui{TOTP Secret} and hit \ui{Save}. From then on, Nexus Terminal answers the ``Verification code:'' prompt itself.
\ui{Test} shows today's code so you can check it matches your phone.

\begin{headsup}
Save the profile first, then save the TOTP secret. And make sure your computer's clock is right: TOTP codes are
time-based, so a clock that's a minute out will produce wrong codes.
\end{headsup}

\section{Do things automatically on login}

\begin{description}
  \item[Startup] A \ui{Startup Command} to run once you're in (say, \code{tmux attach}), plus
        \ui{Environment Variables} as \code{KEY=VALUE}, one per line.
  \item[Login Script] For old-school boxes that ask questions. Add steps: \ui{Wait for} some text, then
        \ui{Send} a reply. Repeat as needed.
\end{description}

\section{Keeping it alive}

Under \ui{Advanced}: the \ui{Terminal Type}, the \ui{Keepalive Interval} and \ui{Max Missed Keepalives}, compression,
and \ui{Auto-Reconnect} with a maximum number of attempts (1–20). Reconnection backs off politely: 1\,s, 2\,s,
4\,s… up to 30\,s.

\begin{sidequest}[What's that number in the status bar?]
For SSH sessions the status bar shows live latency, the round trip to the server. Green is snappy. If it's
in the hundreds and your typing feels sticky, it's the network, not you.
\end{sidequest}

% =========================================================================================
\chapter{Serial: cables and consoles}\label{ch:serial}
\chapmeta{4 minutes}{You'll be talking to a switch, router or microcontroller over USB}

\begin{tldr}
Plug in the adapter, pick the device, match the baud rate. 115200 8N1 is the default and usually right.
Right-click the tab to send a break or wiggle DTR/RTS.
\end{tldr}

\begin{glance}{accent}{Serial}
  Device       & \code{/dev/ttyUSB0}, \code{/dev/cu.usbserial-…} or \code{COM3}\\
  Defaults     & 115200 baud, 8 data bits, no parity, 1 stop bit, no flow control, CR line endings, DTR and RTS on\\
  Extras       & Send Break, toggle DTR/RTS, auto-reconnect when the cable comes back\\
  Good for     & Network gear consoles, Arduinos, ESP32s, Raspberry Pi UARTs, 3D printers\\
\end{glance}

\crop{editor-serial.png}{1515}{40}{0}{70}{0.52\linewidth}{A serial profile for a switch console at 9600 baud. Click refresh after plugging in an adapter.}

\section{Setting it up}

\begin{description}
  \item[Device] Pick it from the list (hit refresh after plugging in) or type the path.
  \item[Line Settings] \ui{Baud Rate} (a preset from 300 to 1\,000\,000, or \ui{Custom}), \ui{Data Bits} (5–8),
        \ui{Parity} (none, odd, even), \ui{Stop Bits} (1 or 2) and \ui{Flow Control} (none, XON/XOFF, RTS/CTS).
  \item[Behaviour] \ui{Line Ending}, \ui{Encoding}, \ui{Local Echo} (for devices that don't echo what you type),
        \ui{Assert DTR on connect}, \ui{Assert RTS on connect} and \ui{Auto-reconnect}.
\end{description}

\section{While you're connected}

Right-click the tab for \ui{Send Break} (handy for getting into a router's ROMMON) and to toggle \ui{DTR} and
\ui{RTS} (handy for resetting boards that wire those to reset). Pull the cable out and the session disconnects; with
auto-reconnect on, it picks up again when you plug it back in.

Serial tabs are proper terminals, so recording, triggers, broadcast, snippets, find and split panes all work.

\begin{headsup}[Gibberish on screen?]
Garbled characters nearly always mean the baud rate, parity or stop bits don't match the device. Check the
device's manual. 9600 and 115200 are the usual suspects.
\end{headsup}

\begin{protip}[Linux and macOS quirks]
On Linux, add yourself to the \code{dialout} group (\code{uucp} on some distributions) and log out and in again.
On a Mac, use the \code{/dev/cu.*} device, not \code{/dev/tty.*}.
\end{protip}

% =========================================================================================
\chapter{Telnet and RAW TCP}\label{ch:raw}
\chapmeta{4 minutes}{You'll poke at old kit and raw sockets without reaching for netcat}

\begin{tldr}
Telnet is a terminal for old devices. RAW is a packet-level conversation with any TCP port, with a hex view.
\end{tldr}

\section{Telnet}

\begin{glance}{amber}{Telnet}
  Default port & 23\\
  Options      & Local echo, line mode, terminal type, encoding (UTF-8, ASCII, ISO 8859-1), line endings, login script\\
  Extras       & Send Break from the tab menu, \ui{Debug Negotiation} for when options go wrong\\
\end{glance}

\begin{danger}[Telnet is plain text]
Everything, including your password, crosses the network unencrypted. Fine on a lab bench, not fine over the
internet or café Wi-Fi.
\end{danger}

\section{RAW TCP}

RAW connects to any TCP port and shows you exactly what goes back and forth. Instead of a terminal you get a
traffic log with \ui{Text}, \ui{Hex} and \ui{Split} views.

\shot{raw.png}{A RAW session to a chatty sensor, with timestamps and direction markers turned on.}

\begin{itemize}
  \item \textbf{Send bar:} pick \ui{Text}, \ui{Hex} or \ui{Escaped} (\code{\textbackslash r\textbackslash n} and friends),
        type, press \kbd{Enter}. One-click buttons send \ui{NUL}, \ui{CR}, \ui{LF}, \ui{ETX} and \ui{EOT}.
  \item \textbf{Presets:} save payloads you send often in the profile under \ui{Payload Presets}, then pick them from
        \ui{Presets…}.
  \item \textbf{Capture:} \ui{Pause}, \ui{Clear}, and \ui{Save} (writes the whole conversation as JSON).
  \item \textbf{Profile options:} encoding (UTF-8, ASCII, Hex), line ending, timestamps, direction, timeouts and
        auto-reconnect.
\end{itemize}

\begin{quickwin}
Point a RAW profile at a web server on port 80 and send \code{HEAD / HTTP/1.0\textbackslash r\textbackslash n\textbackslash r\textbackslash n}
in \ui{Escaped} mode. Instant HTTP headers, no browser required.
\end{quickwin}

% =========================================================================================
\chapter{VNC remote desktops}\label{ch:vnc}
\chapmeta{4 minutes}{You'll see and drive a remote desktop, safely}

\begin{tldr}
VNC shows you another computer's screen. Tunnel it through SSH (tick a box), because plain VNC isn't encrypted.
\end{tldr}

\begin{glance}{rose}{VNC}
  Default port & 5900 (plus the display number: display \code{:1} is 5901)\\
  Display      & \ui{View Only}, \ui{Scale to Fit}, \ui{Shared Session}, \ui{Quality (0-9)}\\
  Security     & Optional SSH tunnel with password, key or agent login\\
\end{glance}

\shot{vnc.png}{A remote desktop, scaled to fit the tab.}

\section{The toolbar}

Connection state and desktop size on the left. On the right: switch between scaled and 1:1, refresh,
\ui{Paste Clipboard} (sends your clipboard to the remote machine) and \ui{Ctrl+Alt+Del}. Click the desktop to give
it your keyboard. Anything the remote machine copies lands on your clipboard automatically.

\section{Tunnel it}

Tick \ui{Enable SSH Tunnel} and fill in the SSH host, port, username and login method. The VNC traffic then
travels inside an encrypted SSH connection. If you connect to a VNC profile \emph{without} a tunnel, Nexus Terminal
warns you and offers \ui{Connect Anyway}, \ui{Edit Profile} or \ui{Cancel}.

\begin{headsup}[Connecting to a Mac?]
Turn on \ui{VNC viewers may control screen with password} in macOS Screen Sharing settings. Nexus Terminal uses
standard VNC password authentication, not Apple's own login method.
\end{headsup}

% =========================================================================================
\chapter{IRC}\label{ch:irc}
\chapmeta{3 minutes}{You'll be chatting in your project's channels from the same window as your servers}

\begin{tldr}
Nick, server, rooms to join. Tabs for each room, \code{/commands} as you'd expect, \kbd{Tab} completes nicks.
\end{tldr}

\begin{glance}{amber}{IRC}
  Default port & 6667, or 6697 with \ui{Use TLS} (use TLS!)\\
  Extras       & NickServ identify on connect, auto-join rooms, room browser, SSH tunnel\\
\end{glance}

\shot{irc.png}{Chatting in \#nexus. Rooms down the right, users in the room below them.}

\begin{itemize}
  \item A \ui{Server} tab plus one per room or private conversation. A dot means unread messages.
  \item Commands: \code{/join}, \code{/part}, \code{/msg}, \code{/me}, \code{/notice}, \code{/nick}, \code{/topic},
        \code{/list}, \code{/names}, \code{/whois}, \code{/who}, \code{/raw}, \code{/quit} and \code{/help}.
  \item \kbd{Tab} completes a nick; click a nick to mention them.
  \item The side panel lists your rooms, a room browser with a filter (click to join) and who's in the current room.
\end{itemize}

% =========================================================================================
\chapter{GitHub}\label{ch:github}
\chapmeta{2 minutes}{You'll see all your repositories, freshest first}

\crop{editor-github.png}{1515}{780}{0}{70}{0.5\linewidth}{A GitHub profile: just the API address and a token.}

\begin{tldr}
Paste a personal access token with the \code{repo} scope. Get a tidy, filterable list of every repository you can see.
\end{tldr}

\shot{github.png}{Repositories grouped by owner, with the latest commit and who made it.}

\begin{itemize}
  \item Grouped by owner or organisation, sorted by most recent activity.
  \item Badges for \proto[rose]{Private}, \proto[violet]{Fork} and \proto[amber]{Archived}, plus language and stars.
  \item Filter box, \ui{Refresh now}, \ui{Visit} (opens in your browser) and \ui{ZIP} (downloads an archive).
  \item Works with GitHub Enterprise: set \ui{API Base URL} to your server's API address. Refreshes every 60–3600 seconds.
\end{itemize}

% =========================================================================================
\chapter{Your own computer}\label{ch:local}
\chapmeta{1 minute}{You'll stop alt-tabbing to a separate terminal app}

Click \ui{Local Machine} for a shell on this computer. It uses your system's default shell unless you choose
another in \menu{Settings\then General\then Local shell} (\code{pwsh}, \code{zsh}, \code{fish}, whatever you fancy).

\shot[0.86\linewidth]{local.png}{A local shell, with all the same terminal powers as SSH.}
